\section{Gang, thép và hợp kim}

\subsection{Hợp kim}

Hợp kim là vật liệu kim loại chứa một kim loại chính và một hay nhiều nguyên tố khác. Tính chất của hợp kim có thể khác đáng kể so với kim loại tinh khiết, vì vậy hợp kim thường được chọn để tăng độ cứng, độ bền, khả năng chống ăn mòn hoặc giảm khối lượng.

\begin{center}
\begin{tabularx}{0.96\textwidth}{>{\bfseries}C{2.7cm}YY}
\toprule
Vật liệu & Thành phần chính & Đặc điểm, ứng dụng \\
\midrule
Gang & Fe, C và một số nguyên tố khác & Cứng, giòn; dùng đúc chi tiết \\
Thép & Fe, C với hàm lượng C thấp hơn gang & Bền, dẻo hơn; dùng trong xây dựng, cơ khí \\
Duralumin & Al cùng một số kim loại khác & Nhẹ, bền; dùng cho chi tiết cần khối lượng nhỏ \\
\bottomrule
\end{tabularx}
\end{center}

\subsection{Từ quặng sắt đến gang}

Trong lò cao, carbon monoxide khử oxide sắt theo nhiều giai đoạn. Có thể biểu diễn khái quát:
\[
\ce{Fe2O3 -> Fe3O4 -> FeO -> Fe}.
\]
Một số phản ứng tiêu biểu:
\[
\ce{3Fe2O3 + CO -> 2Fe3O4 + CO2},
\]
\[
\ce{Fe3O4 + CO -> 3FeO + CO2},
\]
\[
\ce{FeO + CO -> Fe + CO2}.
\]
Đá vôi được dùng để tạo chất trợ dung:
\[
\ce{CaCO3 ->[t^\circ] CaO + CO2},
\]
\[
\ce{CaO + SiO2 -> CaSiO3}.
\]

\subsection{Từ gang đến thép}

Gang chứa nhiều carbon hơn thép. Khi luyện thép, một phần carbon và tạp chất được oxi hóa để giảm hàm lượng của chúng. Ở mức THCS, cần hiểu mục đích chính của quá trình là điều chỉnh thành phần để thu vật liệu có tính chất phù hợp hơn.

\subsection{Kim loại và phi kim}

\begin{center}
\begin{tabularx}{0.96\textwidth}{>{\bfseries}C{3cm}YY}
\toprule
Tiêu chí & Kim loại & Phi kim \\
\midrule
Dẫn điện, dẫn nhiệt & Thường tốt & Thường kém \\
Tính dẻo & Thường có & Thường không có \\
Ion thường tạo & Cation & Anion \\
Oxide thường gặp & Nhiều oxide có tính base & Nhiều oxide có tính acid \\
\bottomrule
\end{tabularx}
\end{center}

\begin{tuluyen}
Vì sao trong nhiều công trình người ta dùng thép thay vì sắt tinh khiết?
\dapso{Thép có cơ tính phù hợp hơn, thường bền và cứng hơn sắt tinh khiết; có thể điều chỉnh thành phần để đáp ứng yêu cầu sử dụng.}
\end{tuluyen}
